\section{Grundlegendes Funktionsprinzip von Wortuhren}
Eine Wortuhr stellt die aktuelle Zeit nicht numerisch, sondern durch hervorgehobene Wörter in natürlicher Sprache dar. 
Hierzu wird ein fest definiertes Wortfeld verwendet, in dem durch gezielte Aktivierung einzelner Buchstaben eine lesbare Zeitangabe entsteht.
Diese sind in einer Matrix angeordnet und bilden verschiedene, teilweise überlappende Wörter. 
Die Darstellung erfolgt durch gezielte Beleuchtung einzelner Buchstaben mittels einer hinterlegten Lichtquelle.

Die Zeit wird durch eine Referenzquelle bereitgestellt und in einer Logikeinheit numerisch verarbeitet. 
Anschließend erfolgt eine regelbasierte Umwandlung in sprachliche Ausdrücke.
Dies wird als im folgenden als Wordmapping bezeichnet. 
Die Zuordnung der Zeitwerte zu den darzustellenden Wörtern ist softwareseitig definiert und wird durch ein geeignetes Logikmodul umgesetzt. 
Dabei kommen typische Formulierungen wie „nach“, „vor“ oder „halb“ zum Einsatz. 
Die Anzeige basiert auf diskreten Zeitintervallen, deren Auflösung die Struktur und Größe der Wortmatrix bestimmt.

Bei der sprachlichen Zeitdarstellung sind mehrere Sonderfälle zu berücksichtigen. 
Hierzu zählt die korrekte Behandlung der Stunde bei „ein Uhr“ im Vergleich zur Verwendung von „eins“ in anderen Kontexten, beispielsweise „halb eins“. 
Weiterhin ist der Übergang zwischen zwei Stunden zu beachten, insbesondere bei Formulierungen wie „vor“ und „halb“, bei denen bereits auf die folgende Stunde Bezug genommen wird.

Zusätzlich erfordern die Zeitpunkte 0:00 Uhr und 12:00 Uhr eine eindeutige Definition. 
Auch Randbereiche zwischen zwei Intervallen, beispielsweise kurz vor oder nach einem Fünf-Minuten-Schritt, müssen eindeutig zugeordnet werden, um eine konsistente Anzeige sicherzustellen.
% Eine Wortuhr stellt die aktuelle Zeit nicht numerisch, sondern durch hervorgehobene Wörter in natürlicher Sprache dar. 
% Hierzu wird ein fest definiertes Wortfeld verwendet, in dem durch gezielte Aktivierung einzelner Buchstaben eine lesbare Zeitangabe entsteht.
% Die Buchstaben sind in einer Matrix angeordnet und bilden verschiedene, teilweise überlappende Wörter. 
% Die Darstellung erfolgt durch Beleuchtung der zugehörigen Buchstaben mittels einer hinterlegten Lichtquelle, welche  je einen Buchstaben beleuchtet.

% Die Zeit muss durch eine entsprechende Referenzquelle zur Verfügung gestellt werden.
% Sie wird mit einer Logikeinheit numerisch verarbeitet und regelbasiert in sprachliche Ausdrücke überführt. 
% Die Regeln, wann welche Buchstaben gesetzt werden, müssen Softwareseitig programmiert sein und auf einem geeigneten Logikmodul prozessiert werden.
% Dabei werden typische Formulierungen wie „nach“, „vor“ oder „halb“ verwendet. 
% Die Anzeige basiert auf diskreten Zeitintervallen, was abhängig der Zeitlichen Genauigkeit die Struktur und Größe der Wortmatrix bestimmt.

% Bei der sprachlichen Zeitdarstellung sind mehrere Sonderfälle zu berücksichtigen. 
% Hierzu zählt die korrekte Behandlung der Stunde bei „ein Uhr“ im Vergleich zur Verwendung von „eins“ in anderen Kontexten, z.B. „Es ist halb eins“. 
% Weiterhin ist der Übergang zwischen zwei Stunden zu beachten, insbesondere bei Formulierungen wie „vor“ und „halb“, bei denen bereits auf die folgende Stunde Bezug genommen wird. 

% Zusätzlich erfordern die Zeitpunkte 0:00 und 12:00 eine eindeutige Definition, da diese sowohl als Beginn eines neuen Tages als auch als Abschluss eines Zeitraums interpretiert werden können. 
% Auch Randbereiche zwischen zwei Intervallen, beispielsweise kurz vor oder nach einem Fünf-Minuten-Schritt, müssen eindeutig zugeordnet werden, um eine konsistente Anzeige sicherzustellen.